\documentclass[../../main.tex]{subfiles}
\begin{document}

\begin{center}
  {\fontsize{48pt}{58pt}\selectfont \textbf{1992}\normalsize\textbf{年度}}
\end{center}

\begin{table}[htb]
  \centering
  \caption{出題テーマと難易度}
  \begin{tabular}{|c|p{9cm}|c|}
    \hline
    \textbf{問題番号} & \textbf{テーマ} & \textbf{難易度} \\
    \hline
    第1問 & 積分で定めた多項式と関数方程式 & \\
    \hline
    第2問 & 立方体の表面の最短経路の数 & \\
    \hline
    第3問 & 放物線と座標軸に接する円 & \\
    \hline
    第4問 & 一次変換と直交ベクトル & \\
    \hline
    第5問 & 玉の交換操作と確率 & \\
    \hline
    第6問 & 正方形を底面とする角錐の体積比 & \\
    \hline
  \end{tabular}
\end{table}

\begin{table}[htb]
  \centering
  \caption{解答}
  \renewcommand{\arraystretch}{1.4}
  \begin{tabular}{|c|c|p{9cm}|}
    \hline
    \textbf{問題番号} & \textbf{小問} & \textbf{答え} \\
    \hline
    第1問 & (1) & $F(x)=-x^2+(a+1)x$ \\
    \cline{2-3}
    & (2) & $f(x)=-2x+\sqrt{2}$ \\
    \hline
    第2問 & (1) & ${}_{3n}C_n$ \\
    \cline{2-3}
    & (2) & $6{}_{3n}C_n-6{}_{2n}C_n$ \\
    \hline
    第3問 & (1) & $1<\dfrac{R}{r}$ \\
    \cline{2-3}
    & (2) & $P\left(\dfrac{\sqrt{2}}{8},\dfrac{1}{32}\right)$ \\
    \hline
    第4問 & (1)，(2) &  \\
    \hline
    第5問 & (1) & $P(Y\le m)=\left(\dfrac{m-2}{N}\right)^n$ \\
    \cline{2-3}
    & (2) & $P(X\le m)=(m-N-1)\left(\dfrac{m-2}{N}\right)^n+(N-m+2)\left(\dfrac{m-1}{N}\right)^n$ \\
    \hline
    第6問 & (1) & 一辺の長さ$2$の正方形（証明問題） \\
    \cline{2-3}
    & (2) & $W_1:W_2=1:\sqrt{5}$ \\
    \hline
  \end{tabular}
\end{table}

\paragraph{試験形態}
6問．試験時間は要確認．

\end{document}
